%% 45장 신뢰 영역 방법
\chapter{신뢰 영역 방법}
\label{ch:trust-region}

\abstract{직선 탐색은 방향을 먼저 정하고 걸음의 길이를 찾는다. 신뢰 영역 방법은 순서를 뒤집어, 이차 모형을 믿을 수 있는 반지름을 먼저 정하고 그 공 안에서 모형을 최소화하는 걸음을 찾는다. 모형이 실제 함수를 얼마나 잘 예측했는지를 보고 반지름을 늘리거나 줄인다. 이 장에서는 먼저 Cauchy 점이 보장하는 모형 감소량을 유도하고, 걸음이 그만큼만 줄여도 기울기가 $0$으로 가는 부분열이 있다는 전역 수렴 정리를 증명한다. 이어서 부분문제의 해를 완전히 특징짓는 Moré--Sorensen의 정리를 하나의 항등식으로 증명하고, 그 해를 Cholesky 분해와 Newton 반복으로 찾는 거의 정확한 방법, 행렬 대신 곱만 쓰는 Steihaug의 공액 기울기 방법을 다룬다. 고유벡터가 기울기와 직교하는 어려운 경우, Newton 반복을 $\|p\|$가 아니라 $1/\|p\|$에 하는 까닭, 그리고 Rosenbrock 함수에서 신뢰 영역이 음의 곡률을 다루는 방식을 수치로 확인한다.}

\section{신뢰 영역의 틀}
\label{sec:tr-framework}

$x_k$에서 이차 모형
\[
m_k(p)=f(x_k)+g_k^{\mathsf T}p+\frac12p^{\mathsf T}B_kp
\]
를 생각하자. $B_k$는 Hessian이나 그 근사이고 양정치일 필요가 없다. 신뢰 영역 방법은 반복마다 \term{신뢰 영역 부분문제}{trust-region subproblem}
\[
\min_p\ m_k(p)\quad\text{조건: }\|p\|\le\Delta_k
\]
를 (근사적으로) 풀어 걸음 $p_k$를 얻는다. $\Delta_k>0$은 모형을 믿을 수 있다고 보는 반지름이다. $B_k$가 양정치가 아니어도 공은 닫혀 있고 유계이므로 부분문제에는 언제나 해가 있고 44장의 Newton 방법처럼 Hessian을 고칠 필요도 없다.

걸음을 받아들일지와 반지름을 어떻게 바꿀지는 실제 감소와 예측 감소의 비
\[
\rho_k=\frac{f(x_k)-f(x_k+p_k)}{m_k(0)-m_k(p_k)}
\]
로 정한다. $\rho_k$가 $1$에 가까우면 모형이 잘 맞았으므로 반지름을 늘려도 되고 작거나 음수이면 모형을 너무 멀리까지 믿었으므로 줄인다. scipy와 stardust의 규칙은 다음과 같다. $\rho_k<1/4$이면 $\Delta_{k+1}=\Delta_k/4$이고, $\rho_k>3/4$이고 걸음이 경계에 닿았으면 $\Delta_{k+1}=\min(2\Delta_k,\Delta_{\max})$이며, 그 밖에는 그대로 둔다. 걸음은 $\rho_k>\eta$일 때만 받아들인다. 기본값은 $\Delta_0=1$, $\Delta_{\max}=1000$, $\eta=0.15$이다. 걸음이 경계 안쪽에 있었다면 반지름이 걸음을 막지 않았으므로 늘릴 까닭이 없다.

\begin{algorithm}[t]
\begin{algorithmic}[1]
\Require 함수 $f$, 기울기, Hessian 또는 그 곱, 시작점 $x_0$, $\Delta_0$, $\Delta_{\max}$, $\eta\in[0,1/4)$, 허용오차 $\mathrm{gtol}$
\Ensure 근사 최소점 $x_k$
\State $k\gets0$
\While{$\|g_k\|\ge\mathrm{gtol}$}
  \State 부분문제 $\min m_k(p)$, $\|p\|\le\Delta_k$를 근사적으로 풀어 $p_k$를 얻는다
  \State $\rho_k\gets\bigl(f(x_k)-f(x_k+p_k)\bigr)/\bigl(m_k(0)-m_k(p_k)\bigr)$
  \If{$\rho_k<1/4$} $\Delta_{k+1}\gets\Delta_k/4$
  \ElsIf{$\rho_k>3/4$이고 $\|p_k\|=\Delta_k$} $\Delta_{k+1}\gets\min(2\Delta_k,\Delta_{\max})$
  \Else\ $\Delta_{k+1}\gets\Delta_k$
  \EndIf
  \If{$\rho_k>\eta$} $x_{k+1}\gets x_k+p_k$ \Else\ $x_{k+1}\gets x_k$ \EndIf
  \State $k\gets k+1$
\EndWhile
\end{algorithmic}
\caption{신뢰 영역 방법}
\label{alg:trust-region}
\end{algorithm}

\section{Cauchy 점과 전역 수렴}
\label{sec:tr-cauchy}

\subsection{Cauchy 점}

부분문제를 정확히 풀지 않아도 되지만, 걸음은 최소한 가장 가파른 내리막 방향에서 얻을 수 있는 만큼은 모형을 줄여야 한다. 그 기준이 되는 \term{Cauchy 점}{Cauchy point}은 $-g_k$ 방향으로 공 안에서 모형을 최소화한 점으로, $g_k\ne0$이면
\[
p_k^{\mathrm C}=-\tau_k\frac{\Delta_k}{\|g_k\|}g_k,\qquad
\tau_k=\begin{cases}1,&g_k^{\mathsf T}B_kg_k\le0,\\\min\bigl(\|g_k\|^3/(\Delta_kg_k^{\mathsf T}B_kg_k),\,1\bigr),&g_k^{\mathsf T}B_kg_k>0\end{cases}
\]
이다. 곡률이 양수가 아니면 모형이 그 방향으로 계속 줄므로 경계까지 가고, 양수이면 한 변수 이차 함수의 최소점이 경계 안에 있는지에 따라 정해진다.

\begin{lemma}
\label{lem:cauchy-decrease}
Cauchy 점은
\[
m_k(0)-m_k(p_k^{\mathrm C})\ge\frac12\|g_k\|\min\Bigl(\Delta_k,\frac{\|g_k\|}{\|B_k\|}\Bigr)
\]
를 만족한다. $\|B_k\|=0$이면 오른쪽의 최솟값은 $\Delta_k$로 읽는다.
\end{lemma}

\begin{proof}
첨자 $k$를 떼고 $\omega=g^{\mathsf T}Bg$로 쓰자. $\omega\le0$이면 $p^{\mathrm C}=-\Delta g/\|g\|$이고
\[
m(0)-m(p^{\mathrm C})=\Delta\|g\|-\frac{\Delta^2}{2\|g\|^2}\omega\ge\Delta\|g\|
\]
이다. $\omega>0$이고 $\|g\|^3\le\Delta\omega$이면 $\tau=\|g\|^3/(\Delta\omega)$이고 $p^{\mathrm C}=-(\|g\|^2/\omega)g$이므로
\[
m(0)-m(p^{\mathrm C})=\frac{\|g\|^4}{\omega}-\frac12\frac{\|g\|^4}{\omega}=\frac{\|g\|^4}{2\omega}\ge\frac{\|g\|^4}{2\|B\|\|g\|^2}=\frac{\|g\|^2}{2\|B\|}
\]
이다. $\omega>0$이고 $\|g\|^3>\Delta\omega$이면 $\tau=1$이고, $\omega<\|g\|^3/\Delta$를 쓰면
\[
m(0)-m(p^{\mathrm C})=\Delta\|g\|-\frac{\Delta^2}{2\|g\|^2}\omega>\Delta\|g\|-\frac{\Delta^2}{2\|g\|^2}\cdot\frac{\|g\|^3}\Delta=\frac12\Delta\|g\|
\]
이다. 세 경우 모두 주장의 오른쪽 이상이다.
\end{proof}

\subsection{전역 수렴}

걸음이 Cauchy 점의 일정 비율만큼 모형을 줄이기만 하면 방법 전체가 수렴한다. 직선 탐색에서 Zoutendijk의 정리가 맡던 역할을 여기서는 다음 정리가 맡는다.

\begin{theorem}
\label{thm:tr-convergence}
$f$가 아래로 유계이고 $\nabla f$가 상수 $L$로 Lipschitz 연속이라 하자. 어떤 상수 $\beta,\nu\ge1$과 $c\in(0,1]$에 대해 모든 $k$에서 $\|B_k\|\le\beta$, $\|p_k\|\le\nu\Delta_k$이고
\[
m_k(0)-m_k(p_k)\ge c\,\|g_k\|\min\Bigl(\Delta_k,\frac{\|g_k\|}{\beta}\Bigr)
\]
이면, \cref{alg:trust-region}의 반복은 $\liminf_{k\to\infty}\|g_k\|=0$을 만족한다.
\end{theorem}

\begin{proof}
그렇지 않다고 하자. 그러면 어떤 $\epsilon>0$과 $K$가 있어 $k\ge K$에서 $\|g_k\|\ge\epsilon$이다. 먼저 반지름이 작으면 $\rho_k$가 $1$에 가깝다는 것을 보인다. 적분 꼴의 Taylor 정리로
\[
f(x_k+p_k)=f(x_k)+g_k^{\mathsf T}p_k+\int_0^1\bigl(\nabla f(x_k+tp_k)-g_k\bigr)^{\mathsf T}p_k\,dt
\]
이므로 $|m_k(p_k)-f(x_k+p_k)|\le\beta\|p_k\|^2/2+L\|p_k\|^2/2\le(\beta+L)\nu^2\Delta_k^2/2$이다. $\Delta_k\le\epsilon/\beta$이면 예측 감소가 $c\epsilon\Delta_k$ 이상이므로
\[
|\rho_k-1|=\frac{|m_k(p_k)-f(x_k+p_k)|}{m_k(0)-m_k(p_k)}\le\frac{(\beta+L)\nu^2\Delta_k}{2c\epsilon}
\]
이다. 그래서 $\bar\Delta=\min\bigl(\epsilon/\beta,\ c\epsilon/(2(\beta+L)\nu^2)\bigr)$로 두면 $\Delta_k\le\bar\Delta$일 때 $\rho_k\ge3/4$이고 반지름이 줄지 않는다. 반지름은 한 번에 $1/4$배까지만 줄므로, $k\ge K$에서 $\Delta_k\ge\min(\Delta_K,\bar\Delta/4)=:\Delta_\flat>0$이다.

이제 $\rho_k\ge1/4$인 $k\ge K$가 무한히 많다고 하자. 그런 $k$에서는 $\rho_k>\eta$이므로 걸음을 받아들이고
\[
f(x_k)-f(x_{k+1})\ge\frac14\bigl(m_k(0)-m_k(p_k)\bigr)\ge\frac c4\,\epsilon\min\Bigl(\Delta_\flat,\frac\epsilon\beta\Bigr)
\]
이다. 함수값은 다른 반복에서 늘지 않으므로 $f(x_k)\to-\infty$가 되어 아래로 유계라는 가정에 어긋난다. 그렇지 않다면 어떤 $K'$부터 언제나 $\rho_k<1/4$이고, 반지름이 매번 $1/4$배가 되어 $\Delta_k\to0$이므로 $\Delta_k\ge\Delta_\flat$에 어긋난다.
\end{proof}

$\eta>0$이면 결론을 $\lim\|g_k\|=0$으로 강화할 수 있다\cite{nocedal2006,conn2000trust}. 이 정리에는 $B_k$가 양정치라는 가정이 없다. Cauchy 점의 감소량만 확보하면 모형이 아무리 나빠도 반지름이 알아서 줄어들어 기울기가 작아진다. 부분문제의 정확한 해는 모형을 Cauchy 점 이상으로 줄이므로 자동으로 이 조건을 만족한다.

\section{부분문제의 해}
\label{sec:tr-subproblem}

\subsection{Moré--Sorensen의 정리}

부분문제는 볼록이 아닐 수도 있는 이차 함수를 공 위에서 최소화하는 문제이다. 일반적인 비볼록 최적화와 달리 이 문제의 전역 최소점은 완전히 특징지어진다.

\begin{theorem}[Moré--Sorensen]
\label{thm:more-sorensen}
$B$가 대칭이고 $\Delta>0$이라 하자. $p_\ast$가 $\min\,g^{\mathsf T}p+p^{\mathsf T}Bp/2$, $\|p\|\le\Delta$의 전역 최소점인 것은 다음을 만족하는 $\lambda\ge0$이 있는 것과 동치이다.
\[
\text{(i) }(B+\lambda I)p_\ast=-g,\qquad\text{(ii) }\lambda(\Delta-\|p_\ast\|)=0,\qquad\text{(iii) }B+\lambda I\text{가 양반정치.}
\]
\end{theorem}

\begin{proof}
상수항을 뺀 모형을 $m(p)=g^{\mathsf T}p+p^{\mathsf T}Bp/2$로 쓰자. (i)을 만족하는 $p_\ast$와 $\lambda$가 있으면 모든 $p$에 대해
\begin{equation}
m(p)-m(p_\ast)=\frac12(p-p_\ast)^{\mathsf T}(B+\lambda I)(p-p_\ast)+\frac\lambda2\bigl(\|p_\ast\|^2-\|p\|^2\bigr)
\label{eq:tr-identity}
\end{equation}
이다. $g=-(B+\lambda I)p_\ast$를 넣고 양변을 펼치면 확인된다.

충분성. (i)--(iii)이 성립하고 $\|p\|\le\Delta$이면 \eqref{eq:tr-identity}의 첫 항은 (iii)으로 음이 아니다. 둘째 항은 $\lambda=0$이면 $0$이고, $\lambda>0$이면 (ii)로 $\|p_\ast\|=\Delta\ge\|p\|$이므로 음이 아니다. 따라서 $m(p)\ge m(p_\ast)$이다.

필요성. $p_\ast$가 전역 최소점이라 하자. $\|p_\ast\|<\Delta$이면 $p_\ast$는 이차 함수 $m$의 제약 없는 국소 최소점이므로 42장의 최적성 조건으로 $Bp_\ast=-g$이고 $B$가 양반정치이다. $\lambda=0$으로 두면 된다. $\|p_\ast\|=\Delta$이면 $p_\ast$는 구면 $\|p\|=\Delta$ 위의 최소점이기도 하므로 Lagrange 승수 정리로 (i)을 만족하는 $\lambda$가 있고, (ii)는 저절로 성립한다. 남은 것은 $\lambda\ge0$과 (iii)이다. $w^{\mathsf T}p_\ast\ne0$인 아무 $w$에 대해 반사한 점 $p=p_\ast-2(w^{\mathsf T}p_\ast/\|w\|^2)w$도 구면 위에 있으므로 \eqref{eq:tr-identity}에서 둘째 항이 사라지고
\[
0\le m(p)-m(p_\ast)=2\Bigl(\frac{w^{\mathsf T}p_\ast}{\|w\|^2}\Bigr)^2w^{\mathsf T}(B+\lambda I)w
\]
이다. 그래서 $w^{\mathsf T}p_\ast\ne0$인 모든 $w$에서 $w^{\mathsf T}(B+\lambda I)w\ge0$이고, 그런 $w$는 조밀하므로 연속성으로 (iii)이 나온다. 끝으로 $\lambda<0$이라 하자. 공 안의 점 $p=(1-s)p_\ast$($0<s<1$)에서 \eqref{eq:tr-identity}는
\[
m(p)-m(p_\ast)=\frac{s^2}2p_\ast^{\mathsf T}(B+\lambda I)p_\ast+\frac\lambda2\Delta^2(2s-s^2)
\]
이고, $s\to0$이면 $\lambda\Delta^2s$ 항이 지배해 음수가 되므로 $p_\ast$가 최소점이라는 데 어긋난다.
\end{proof}

이 정리\cite{more1983trust,gay1981,sorensen1982trust}에 따르면 부분문제는 $\lambda\ge\max(0,-\lambda_1)$을 찾는 문제가 된다. 여기서 $\lambda_1$은 $B$의 가장 작은 고유값이다. $B$가 양정치이고 Newton 걸음 $-B^{-1}g$가 공 안에 있으면 $\lambda=0$이 답이다. 그렇지 않으면 경계 위의 해, 곧
\[
\|p(\lambda)\|=\Delta,\qquad p(\lambda)=-(B+\lambda I)^{-1}g
\]
를 만족하는 $\lambda>\max(0,-\lambda_1)$을 찾는다. 이 식을 \term{영년 방정식}{secular equation}이라 한다. $\lambda$는 47장에서 볼 Levenberg--Marquardt 감쇠와 같은 역할을 하며, $\lambda$가 커질수록 걸음은 짧아지고 가장 가파른 내리막 방향으로 돈다.

\subsection{영년 방정식과 Newton 반복}

$B=Q\Lambda Q^{\mathsf T}$로 대각화하고 $\gamma_j=q_j^{\mathsf T}g$로 쓰면
\[
\|p(\lambda)\|^2=\sum_{j=1}^n\frac{\gamma_j^2}{(\lambda_j+\lambda)^2}
\]
이다. $\gamma_1\ne0$이면 $\lambda\downarrow-\lambda_1$에서 $\|p\|\to\infty$이고 $\lambda\to\infty$에서 $\|p\|\to0$이며 그 사이에서 순감소하므로 해가 하나 있다. 이 함수에 그대로 Newton 방법을 쓰면, $\|p(\lambda)\|$가 극점 근처에서 $|\gamma_1|/(\lambda+\lambda_1)$처럼 휘어 있어 접선이 멀리 빗나가므로 느리다. 대신 역수
\[
\phi(\lambda)=\frac1{\|p(\lambda)\|}-\frac1\Delta
\]
에 Newton 방법을 쓴다. 극점 하나만 있으면 $1/\|p(\lambda)\|=(\lambda_1+\lambda)/|\gamma_1|$이 정확히 일차 함수이므로 Newton 한 번으로 끝나고, 일반적으로도 거의 일차 함수이다. 도함수는 Cholesky 분해 $B+\lambda I=R^{\mathsf T}R$로 싸게 계산된다. $p$를 $R^{\mathsf T}Rp=-g$로 풀고 $q$를 $R^{\mathsf T}q=p$로 풀면 $d\|p\|/d\lambda=-p^{\mathsf T}(B+\lambda I)^{-1}p/\|p\|=-\|q\|^2/\|p\|$이므로
\[
\lambda\gets\lambda+\Bigl(\frac{\|p\|}{\|q\|}\Bigr)^2\frac{\|p\|-\Delta}\Delta
\]
이다. $B=\operatorname{diag}(1,2,5)$, $g=(1,1,1)$, $\Delta=0.1$에서 $\lambda=0$으로 시작하면 해 $\lambda_\ast=14.8846$까지의 오차가 $\phi$의 Newton 반복으로는 $3.1$, $8.8\times10^{-3}$, $5.0\times10^{-8}$, $1.8\times10^{-15}$으로 네 번에 끝나지만, $\|p\|-\Delta$의 Newton 반복으로는 $14$, $12$, $8.6$, $4.4$, $1.2$, $7.9\times10^{-2}$, \dots로 아홉 번이 든다.

\subsection{어려운 경우}

$\lambda_1\le0$이고 $g$가 $\lambda_1$의 고유공간과 직교하면(고유값이 단순하면 $\gamma_1=q_1^{\mathsf T}g=0$이면) $\lambda\downarrow-\lambda_1$에서도 $\|p(\lambda)\|$가 유한한 값에 머문다. 그 값이 $\Delta$보다 작으면 영년 방정식에 해가 없고, 이것을 \term{어려운 경우}{hard case}라 한다. 이때 해는 $\lambda=-\lambda_1$이고
\[
p_\ast=-(B-\lambda_1I)^+g+\tau q_1
\]
에서 $\|p_\ast\|=\Delta$가 되도록 $\tau$를 고른다. $(B-\lambda_1I)^+$는 의사역행렬이다. \cref{thm:more-sorensen}의 세 조건이 모두 성립하는 것은 바로 확인된다. 예컨대 $B=\operatorname{diag}(-1,1)$, $g=(0,1)$, $\Delta=2$이면 $\lambda>1$에서 $p(\lambda)=(0,-1/(1+\lambda))$의 길이가 $1/2$보다 작아 $\Delta$에 닿지 않는다. 해는 $\lambda=1$, $p_\ast=(\pm\sqrt{15}/2,\,-1/2)$로, 음의 곡률 방향 $q_1=e_1$을 따라 경계까지 간다. 정확히 $\gamma_1=0$인 일은 드물지만 $\gamma_1$이 아주 작으면 같은 수치적 어려움이 생기므로 실제 알고리즘은 이 경우를 따로 다룬다.

\subsection{거의 정확한 방법}

Moré와 Sorensen의 알고리즘\cite{more1983trust}은 영년 방정식의 Newton 반복에 안전 장치를 단 것이다. $\lambda$를 낀 구간 $[\lambda_L,\lambda_U]$를 유지하고, $B+\lambda I$의 Cholesky 분해가 실패하면 $\lambda$가 너무 작다는 뜻이므로 $\lambda_L$을 올린다. $\|p(\lambda)\|<\Delta$이면 어려운 경우일 수 있으므로 분해된 $R$에서 $\|Rz\|$가 작은 단위 벡터 $z$를 Cline, Moler, Stewart, Wilkinson의 방법\cite{cline1979}으로 어림하고 경계까지 $z$ 방향으로 간다. $z$는 $B+\lambda I$의 가장 작은 고유값의 고유벡터에 가까운 방향이다. 초기 구간은 Gershgorin 원판과 Frobenius·무한대 노름으로 잡는다. 정확히 풀지 않고
\[
\bigl|\|p\|-\Delta\bigr|\le k_{\mathrm{easy}}\Delta,\qquad\text{또는 어려운 경우에}\quad\tau^2\|Rz\|^2\le k_{\mathrm{hard}}\bigl(p^{\mathsf T}(B+\lambda I)p+\lambda\Delta^2\bigr)
\]
이면 멈춘다. scipy와 stardust의 기본값은 $k_{\mathrm{easy}}=0.1$, $k_{\mathrm{hard}}=0.2$이다. 반지름 자체가 어림이므로 경계를 정확히 맞추려고 분해를 더 해도 얻을 것이 없고 그래서 경계 조건을 $10\%$까지 느슨하게 두어도 된다. 반복마다 분해가 몇 번 들므로 이 방법은 $n$이 수백 이하인 문제에 알맞다.

\section{공액 기울기 신뢰 영역}
\label{sec:steihaug}

$n$이 크면 분해가 너무 비싸다. Steihaug와 Toint는 27장의 공액 기울기법을 $Bp=-g$에 $p_0=0$부터 쓰되, 두 경우에 경계에서 멈추게 했다\cite{steihaug1983,toint1981}. 곡률 $d_j^{\mathsf T}Bd_j\le0$인 방향을 만나면 그 방향으로 경계까지 가고(두 교점 가운데 모형이 작은 쪽), 다음 반복점이 공을 벗어나면 그 방향과 경계의 교점에서 멈춘다. 그 밖에는 44장의 Newton-CG처럼 잔차가 $\min(1/2,\sqrt{\|g\|})\|g\|$ 아래로 내려가면 멈춘다.

이렇게 경계에서 멈추어도 되는 것은 다음 성질 덕분이다. 첫 반복점은 $-g$ 방향이므로 Cauchy 점과 같은 방향에서 모형을 줄이고, 그 뒤로는 모형이 계속 줄면서 반복점의 길이는 계속 는다. 그래서 한 번 공을 벗어난 반복은 다시 돌아오지 않으므로, 경계에서 멈추어도 버리게 되는 것은 공 밖의 반복점뿐이다.

\begin{proposition}[Steihaug]
\label{prop:steihaug}
$p_0=0$에서 시작한 공액 기울기 반복이 $d_0^{\mathsf T}Bd_0,\dots,d_{j-1}^{\mathsf T}Bd_{j-1}>0$인 동안 만드는 반복점은 $0=\|p_0\|<\|p_1\|<\dots<\|p_j\|$를 만족한다.
\end{proposition}

\begin{proof}
공액 기울기법의 기호로 $p_{i+1}=p_i+\alpha_id_i$, $r_{i+1}=r_i+\alpha_iBd_i$, $d_{i+1}=-r_{i+1}+\beta_{i+1}d_i$이고, $\alpha_i=r_i^{\mathsf T}r_i/(d_i^{\mathsf T}Bd_i)>0$, $\beta_{i+1}=r_{i+1}^{\mathsf T}r_{i+1}/(r_i^{\mathsf T}r_i)>0$이다. 27장에서 보았듯이 $r_i$는 $d_0,\dots,d_{i-1}$과 직교하고, $p_i$는 이들의 일차 결합이므로 $p_i^{\mathsf T}r_i=0$이다. 그러면
\[
p_{i+1}^{\mathsf T}d_{i+1}=p_{i+1}^{\mathsf T}(-r_{i+1}+\beta_{i+1}d_i)=\beta_{i+1}(p_i+\alpha_id_i)^{\mathsf T}d_i=\beta_{i+1}\bigl(p_i^{\mathsf T}d_i+\alpha_i\|d_i\|^2\bigr)
\]
이고, $p_0^{\mathsf T}d_0=0$에서 시작하는 귀납법으로 $i\ge1$에서 $p_i^{\mathsf T}d_i>0$이다. 따라서
\[
\|p_{i+1}\|^2=\|p_i\|^2+2\alpha_ip_i^{\mathsf T}d_i+\alpha_i^2\|d_i\|^2>\|p_i\|^2
\]
이다.
\end{proof}

모형이 계속 주는 것은, 27장에서 보았듯이 공액 기울기 반복이 크릴로프 부분공간에서 모형을 최소화하고 그 부분공간이 반복마다 커져 최솟값이 늘 수 없기 때문이다. 그래서 Steihaug의 걸음은 Cauchy 점 이상으로 모형을 줄여 \cref{thm:tr-convergence}의 조건을 만족하고, 공 안에서 끝나면 비정확 Newton 걸음이므로 44장의 초선형 수렴도 이어받는다. 음의 곡률을 만난 경우가 44장의 Newton-CG와 다르다. Newton-CG는 그 방향을 버리고 그때까지의 해나 가장 가파른 내리막 방향을 썼지만, 신뢰 영역에서는 그 방향으로 경계까지 가는 것이 모형을 가장 많이 줄이는 선택이다.

\begin{aside}[Powell의 혼성법과 dogleg]
신뢰 영역의 생각은 Levenberg(1944)와 Marquardt(1963)의 비선형 최소제곱 방법\cite{levenberg1944,marquardt1963}에 이미 있었지만, 일반적인 최적화의 틀로 세운 것은 Powell이다. 그의 1970년 혼성법\cite{powell1970hybrid}은 34장의 MINPACK \code{hybrd}의 바탕이며, 부분문제를 $0$에서 $-g$ 방향의 모형 최소점 $-(g^{\mathsf T}g/g^{\mathsf T}Bg)\,g$까지, 거기서 Newton 걸음 $-B^{-1}g$까지 꺾인 두 선분 위에서 푼다. 그 모양이 개의 뒷다리를 닮아 dogleg이라 불린다. $B$가 양정치이면 이 경로를 따라 모형은 줄고 길이는 늘므로 경계와 많아야 한 번 만난다. scipy의 \code{minimize}에는 \code{dogleg} 방법도 있는데 $B$가 양정치여야 한다. stardust의 \code{minimize}에는 \code{'trust-ncg'}와 \code{'trust-exact'}만 있고, dogleg 걸음은 34장의 \code{root(method='hybr')}가 쓴다.
\end{aside}

\section{수치 비교}
\label{sec:tr-numerics}

2차원 Rosenbrock 함수를 $(-1.2,1)$에서 시작하면 scipy 1.18.1의 기본 설정($\|g\|<10^{-4}$에서 멈춤)으로 trust-exact가 $25$번, trust-ncg가 $29$번 반복한다. 44장에서 같은 문제의 Newton-CG는 $83$번이었다. 차이는 처음 몇 걸음에서 생긴다. Newton-CG는 첫 반복에서 걸음 길이를 제한하지 않아 Hessian이 양정치가 아닌 영역으로 들어가고, 그 영역의 $42$개 반복점에서 음의 곡률을 버리고 가장 가파른 내리막에 가까운 방향으로 기어간다. 신뢰 영역은 $\Delta_0=1$로 첫 걸음을 묶어 두므로 trust-exact의 반복점은 한 번도 그 영역에 들어가지 않는다.

100차원 Rosenbrock 함수(scipy의 \code{rosen})를 $(-1.2,1,-1.2,1,\dots)$에서 시작하면 trust-ncg는 $443$번 반복하고 Hessian과 벡터의 곱을 $2{,}305$번 계산해 $f=1.1\times10^{-12}$인 전역 최소점에 이른다. trust-exact는 $205$번 반복하지만 $f=3.98662$, $x_1\approx-0.993$인 다른 국소 최소점에 멈춘다. 이 함수는 $n=100$에서 전역 최소점 말고도 그런 국소 최소점을 가지며 그곳의 Hessian은 양정치이므로(가장 작은 고유값 $0.499$) 어느 방법이 거기서 멈추어도 이상할 것이 없다. 어느 국소 최소점에 이를지는 경로가 정하고, 경로는 부분문제를 얼마나 정확히 푸는지에 따라 달라진다. 42장에서 말했듯이 볼록이 아닌 함수에서 알고리즘이 약속하는 것은 국소 최소이다.

\begin{instardust}[\code{docs/design/optimize.md} §3.5]
stardust \code{minimize}의 두 신뢰 영역 방법 \code{'trust-ncg'}, \code{'trust-exact'}는 scipy 1.18.1의 \code{\_minimize\_trust\_region}과 두 부분문제(Steihaug--Toint의 공액 기울기, Moré--Sorensen의 반복)를 같은 연산 순서로 옮긴 것으로(\code{csrc/optimize/trust.cpp}), 성공 경로에서는 함수를 부르는 점이 scipy와 비트까지 같고 다른 것은 실패 경로뿐이다. scipy의 공액 기울기 부분문제에는 반복 상한이 없어 대칭이 아닌 Hessian에서 끝나지 않을 수 있으므로 stardust는 $20n$번에서 멈추고 status $3$을 돌려준다. trust-exact의 시험점에서 Hessian이 유한하지 않으면 scipy는 예외를 내지만, stardust는 그 걸음을 $\rho=-\infty$로 보고 거절해 반지름을 줄인다. $\rho$가 NaN일 때도 반지름을 $1/4$배로 줄인다. 분해는 LAPACK \code{potrf}이고, 어려운 경우의 $z$는 Cline--Moler--Stewart--Wilkinson의 어림이다. 두 방법 모두 기울기가 있어야 하며, trust-exact는 행렬을 분해해야 하므로 차분 Hessian을 받지 않는다.
\end{instardust}

\section{맺음}

신뢰 영역 방법은 걸음의 크기를 모형의 신뢰도에 묶는다. 비 $\rho_k$가 그 신뢰도를 재고 반지름이 그에 따라 움직인다. Cauchy 점은 걸음이 모형을 적어도 얼마나 줄여야 하는지를 정하고 그만큼만 줄이면 $B_k$가 양정치가 아니어도 수렴한다. 부분문제 자체는 Moré--Sorensen의 정리 덕분에 완전히 풀리는 문제이고, 그 해는 $\lambda$로 감쇠한 Newton 걸음이다. 작은 문제에서는 Cholesky 분해와 영년 방정식으로 거의 정확히 풀고, 큰 문제에서는 공액 기울기 반복을 경계에서 끊는다. 다음 장에서는 기울기마저 쓰지 않는 방법으로 넘어가고, 47장에서는 같은 신뢰 영역의 틀이 비선형 최소제곱에서 Levenberg--Marquardt와 반사 신뢰 영역으로 다시 나타난다.

신뢰 영역 방법의 모든 것은 Conn, Gould, Toint의 책\cite{conn2000trust}에 있고, Nocedal과 Wright의 책\cite{nocedal2006} 4장과 7장이 좋은 입문이다.

\chapterreferences
